% Introducción y Motivación
\section{Introducción}

\begin{frame}
\frametitle{Contexto y motivación}
\begin{itemize}
    \item \textbf{Tendencia global}: más de 40 países han legalizado el cannabis en alguna forma.
    \item \textbf{Colombia}: el debate sobre regulación incorpora hoy preguntas económicas concretas (tamaño de mercado, base tributaria, política preventiva).
    \item \textbf{Pregunta previa}: antes de cualquier simulación fiscal, hay que caracterizar empíricamente \textbf{quién consume} y \textbf{con qué propensión}.
    \item \textbf{Insumo natural}: la Encuesta Nacional de Consumo de Sustancias Psicoactivas (ENCSPA).
\end{itemize}
\end{frame}

\begin{frame}
\frametitle{Brecha que aborda este trabajo}
\begin{itemize}
    \item Pocos modelos comparan rigurosamente \textbf{benchmark econométrico vs.\ \textit{machine learning}} en este dominio para Colombia.
    \item Escasa discusión estructurada de \textbf{ética y sesgos} en estudios de consumo de drogas con datos de encuesta.
    \item Necesidad de \textbf{interpretabilidad} (no sólo precisión) cuando el modelo se usa para política pública.
\end{itemize}
\vspace{0.3cm}
\textbf{Nuestra contribución}: pipeline reproducible + comparación honesta + interpretabilidad SHAP + frontera ética explícita.
\end{frame}

\begin{frame}
\frametitle{¿Por qué importa?}
\begin{columns}
\column{0.5\textwidth}
\textbf{Dimensión económica:}
\begin{itemize}
    \item Tamaño potencial del mercado
    \item Base tributaria si se legaliza
    \item Elasticidad-precio del segmento
\end{itemize}

\column{0.5\textwidth}
\textbf{Dimensión social:}
\begin{itemize}
    \item Focalización de prevención
    \item Reducción del estigma vía datos
    \item Uso responsable de IA en política
\end{itemize}
\end{columns}
\end{frame}
